%%%%%%%%%%%%%%%%%%%%%%%%%%%%%%%%%%%%%%%%%%%%%%%%%%%%%%%%%%%%%%%%%%%%%%%%%%%%%%
%  config.tex — DATOS DE TU TESIS
%
%  Este es el ÚNICO archivo que necesitas editar para personalizar la portada
%  y los metadatos. No toques tesisbima.cls salvo que sepas lo que haces.
%%%%%%%%%%%%%%%%%%%%%%%%%%%%%%%%%%%%%%%%%%%%%%%%%%%%%%%%%%%%%%%%%%%%%%%%%%%%%%

%------------------------------------------------------------------------------
% TÍTULO Y AUTOR
%   La portada los imprime en MAYÚSCULA y negrita automáticamente:
%   escríbelos normal, con sus tildes.
%------------------------------------------------------------------------------
\titulotesis{Título de la tesis en mayúscula y negrita,
             letra Times Roman 15--16, centrado}

\autor{Nombre Completo del Estudiante}

%------------------------------------------------------------------------------
% PROFESOR PATROCINANTE
%------------------------------------------------------------------------------
\patrocinante{Dr. Nombre Apellido}
\institutopatrocinante{Instituto de Ciencias Marinas y Limnológicas}
\facultadpatrocinante{Facultad de Ciencias -- Universidad Austral de Chile}

%------------------------------------------------------------------------------
% PROFESOR CO-PATROCINANTE (si corresponde)
%   Si no corresponde, deja los tres comandos con las llaves vacías:
%   el bloque completo desaparece de la portada.
%------------------------------------------------------------------------------
\copatrocinante{}
\institutocopatrocinante{}
\facultadcopatrocinante{}

% Ejemplo con co-patrocinante:
% \copatrocinante{Dra. Nombre Apellido}
% \institutocopatrocinante{Instituto de Ciencias Marinas y Limnológicas}
% \facultadcopatrocinante{Facultad de Ciencias -- Universidad Austral de Chile}

%------------------------------------------------------------------------------
% CIUDAD Y AÑO
%------------------------------------------------------------------------------
\ciudad{Valdivia -- Chile}
\anio{2026}

%------------------------------------------------------------------------------
% LOGOS DE LA PORTADA
%
%   El Anexo 3 del Reglamento muestra la portada SIN logos. Si tu comisión
%   los pide (como en el formato del Proyecto de Tesis), descomenta el bloque
%   de abajo. Para volver a la portada oficial exacta, deja \logosportada{}.
%------------------------------------------------------------------------------
\logosportada{}

% --- Portada con logos (escudo UACh + logo del laboratorio) ---
% \logosportada{%
%   \includegraphics[height=2.2cm]{logos/escudo_uach.png}\hspace{1.2cm}%
%   \includegraphics[height=2.2cm]{logos/logo_lab.png}%
%   \\[0.5em]
%   \normalsize Universidad Austral de Chile\\
%   Facultad de Ciencias\\
%   Escuela de Biología Marina
% }
%
% En logos/ también está logo_uach_tiny_name.png (logo horizontal de la UACh
% con el nombre), por si prefieres ese en vez del escudo.
%
% NOTA: en lab-recursos/logos no hay un logo del ICML. Cuando lo consigas,
% cópialo a la carpeta logos/ de esta plantilla y reemplaza logo_lab.png.

%------------------------------------------------------------------------------
% TEXTO DEL GRADO
%   Es el que exige el Anexo 3; cámbialo solo si tu Escuela lo indica.
%------------------------------------------------------------------------------
% \textogrado{Tesis de Grado presentada como parte de los requisitos para
%   optar al grado de \textbf{Licenciado en Biología Marina} y
%   \textbf{Título Profesional de Biólogo Marino.}}

%------------------------------------------------------------------------------
% METADATOS DEL PDF
%------------------------------------------------------------------------------
\hypersetup{
  pdftitle  = {Título de la tesis},
  pdfauthor = {Nombre Completo del Estudiante},
  pdfsubject= {Tesis de Grado — Escuela de Biología Marina, UACh},
  pdfkeywords = {biología marina, UACh, tesis}
}

%------------------------------------------------------------------------------
% RUTAS DE FIGURAS Y ARCHIVO DE BIBLIOGRAFÍA
%------------------------------------------------------------------------------
\graphicspath{{figuras/}{logos/}}
\addbibresource{referencias.bib}

%------------------------------------------------------------------------------
% AJUSTES OPCIONALES
%------------------------------------------------------------------------------
% Párrafos sin sangría y con espacio entre ellos (estilo Word):
% \setlength{\parindent}{0pt}
% \setlength{\parskip}{\baselineskip}

% Mostrar hasta subsubsecciones (x.y.z) en el Índice General:
% \setcounter{tocdepth}{3}
